\begin{table}[t]
\centering
\fontsize{8}{10.5}\selectfont
\setlength{\tabcolsep}{5pt}
\begin{tabular}{l cccc}
\toprule
\textbf{Meta-Reviewer} & \textsc{All} & \textsc{Human} & \textsc{AI} & \textsc{Self} \\
\midrule
\textsc{Random Baseline} & 10.0\% & 10.0\% & 10.0\% & 10.0\% \\
\textsc{Human (agreement)} & 49.3\% & 47.2\% & 52.9\% & --- \\
\midrule
\textsc{Claude-Opus-4.7} & 44.3\% & 40.0\% & 51.2\% & 47.9\% \\
\textsc{GPT-5.4} & 30.1\% & 31.2\% & 28.2\% & 29.1\% \\
\textsc{Gemini-3.1-Pro} & 29.0\% & 26.1\% & 33.8\% & 39.6\% \\
\bottomrule
\end{tabular}
\caption{\textbf{Ten-class meta-reviewer accuracy on the 908-item benchmark (vs.\ primary annotator).}
The 10-class label encodes both the annotation cascade outcome (correctness $\to$ significance $\to$ evidence)
and inter-annotator agreement/disagreement, making it a much harder prediction target than per-axis accuracy.
Columns show accuracy across four review subsets: all items, human review items only, AI review items only,
and ``self review'' where the meta-reviewer judges items from its own model family.
The human baseline (secondary annotator vs.\ primary) achieves 49.3\%, reflecting the substantial
inter-annotator disagreement on the fine-grained 10-class taxonomy.
``Self review'' is not applicable (---) for the human baseline.}
\label{tab:metareview-uncertainty-accuracy}
\end{table}
